\textbf{Construcción de la matriz y producto por filas.} Se convierte cada ángulo a radianes y se escribe $c_\phi=\cos\phi$, $s_\phi=\sin\phi$, y análogamente para los otros ángulos. Multiplicar las rotaciones pasivas en orden 3--2--1 da:\[\fitmath{C_{bn}=\begin{bmatrix}c_\theta c_\psi&c_\theta s_\psi&-s_\theta\\s_\phi s_\theta c_\psi-c_\phi s_\psi&s_\phi s_\theta s_\psi+c_\phi c_\psi&s_\phi c_\theta\\c_\phi s_\theta c_\psi+s_\phi s_\psi&c_\phi s_\theta s_\psi-s_\phi c_\psi&c_\phi c_\theta\end{bmatrix}}\]\[\fitmath{u=(0.901342938)(82)+(0.420303115)(-4)+(-0.104528463)(-3)\approx72.54249\,\mathrm{m/s}}\]\[\fitmath{v=(-0.393686533)(82)+(0.895687419)(-4)+(0.206772729)(-3)\approx-36.48536\,\mathrm{m/s}}\]\[\fitmath{w=(0.180532051)(82)+(-0.145221691)(-4)+(0.972789206)(-3)\approx12.46615\,\mathrm{m/s}}\]La inversa es la transpuesta: multiplicar $C_{bn}^T$ por el vector obtenido recupera las tres componentes NED. Las normas coinciden: $V=82.1523\,\mathrm{m/s}$. Para incidencia se usa $\alpha=\operatorname{atan2}(w,u)$; el argumento de $w/u$ por sí solo pierde el cuadrante.
